\section{AMC 10 10B 2011}
        \begin{enumerate}
        	\item (\textit{AMC 10 10B 2011}) What is 
 \[\dfrac{2+4+6}{1+3+5} - \dfrac{1+3+5}{2+4+6}\]


 $\textbf{(A)}\ -1\qquad\textbf{(B)}\ \frac{5}{36}\qquad\textbf{(C)}\ \frac{7}{12}\qquad\textbf{(D)}\ \frac{147}{60}\qquad\textbf{(E)}\ \frac{43}{3}$


 
 \textbf{Answer: }C



\textbf{Solution}

$\dfrac{2+4+6}{1+3+5} - \dfrac{1+3+5}{2+4+6} = \dfrac{12}{9} - \dfrac{9}{12} = \dfrac{4}{3} - \dfrac{3}{4} = \boxed{\dfrac{7}{12}\; \textbf{(C)}}$

 
Note: This exact problem was reused in 2013 AMC 10B: 

 \href{https://artofproblemsolving.com/wiki/index.php/2013\_AMC\_10B\_Problems/Problem\_1}{https://artofproblemsolving.com/wiki/index.php/2013\_AMC\_10B\_Problems/Problem\_1}

 


	\item (\textit{AMC 10 10B 2011}) Josanna's test scores to date are $90, 80, 70, 60,$ and $85$. Her goal is to raise her test average at least $3$ points with her next test. What is the minimum test score she would need to accomplish this goal?


 $\textbf{(A)}\ 80 \qquad\textbf{(B)}\ 82 \qquad\textbf{(C)}\ 85 \qquad\textbf{(D)}\ 90 \qquad\textbf{(E)}\ 95$


 
 \textbf{Answer: }E



\textbf{Solution}

The average of her current scores is $77$. To raise it $3$ points, she needs an average of $80$, and so after her $6$ tests, a sum of $480$. Her current sum is $385$, so she needs a $480 - 385 = \boxed{\mathrm{(E) \ } 95}$

 


	\item (\textit{AMC 10 10B 2011}) At a store, when a length is reported as $x$ inches that means the length is at least $x - 0.5$ inches and at most $x + 0.5$ inches. Suppose the dimensions of a rectangular tile are reported as $2$ inches by $3$ inches. In square inches, what is the minimum area for the rectangle?


 $\textbf{(A)}\ 3.75 \qquad\textbf{(B)}\ 4.5 \qquad\textbf{(C)}\ 5 \qquad\textbf{(D)}\ 6 \qquad\textbf{(E)}\ 8.75$


 
 \textbf{Answer: }A



\textbf{Solution 1}

The minimum dimensions of the rectangle are $1.5$ inches by $2.5$ inches. The minimum area is $1.5\times2.5=\boxed{\mathrm{(A) \ } 3.75}$ square inches.

 



\textbf{Solution 2 (Utilize algebra)}

Notice that the minimum area of a rectangle with dimensions $x$ inches by $y$ inches will result in $(x - 0.5)\times(y - 0.5)$. Simplify to get $xy - 0.5x - 0.5y + 0.25$. Plug in $2$ and $3$ for $x$ and $y$ to get $6 - 1 - 1.5 + 0.25 = \boxed{\text{(A)\ 3.75}}$ square inches.

 


	\item (\textit{AMC 10 10B 2011}) LeRoy and Bernardo went on a week-long trip together and agreed to share the costs equally. Over the week, each of them paid for various joint expenses such as gasoline and car rental. At the end of the trip, it turned out that LeRoy had paid $A$ dollars and Bernardo had paid $B$ dollars, where $A < B$. How many dollars must LeRoy give to Bernardo so that they share the costs equally?


 $\textbf{(A)}\ \frac{A + B}{2} \qquad\textbf{(B)}\ \dfrac{A - B}{2}\qquad\textbf{(C)}\ \dfrac{B - A}{2}\qquad\textbf{(D)}\ B - A \qquad\textbf{(E)}\ A + B$


 
 \textbf{Answer: }C



\textbf{Solution 1}

The difference in how much LeRoy and Bernardo paid is $B-A$. To share the costs equally, LeRoy must give Bernardo half of the difference, which is $\boxed{\textbf{(C) } \;\frac{B-A}{2}}$

 



\textbf{Solution 2 (Use an example)}

Since there are no restrictions on cost paid besides $A<B$, we can use an example where $A = 40$ and $B = 50$. Quickly, we realize the only way they could pay the same amount of money is if they both pay 45 dollars. This means LeRoy must give Bernardo $50 - 45 = 5$. Looking at the answer choices we see only $\frac{50-40}{2} = 5$ works. $\boxed{\textbf{(C) } \;\frac{B-A}{2}}$

 


	\item (\textit{AMC 10 10B 2011}) In multiplying two positive integers $a$ and $b$, Ron reversed the digits of the two-digit number $a$. His erroneous product was $161$. What is the correct value of the product of $a$ and $b$?


 $\textbf{(A)}\ 116 \qquad\textbf{(B)}\  161 \qquad\textbf{(C)}\  204 \qquad\textbf{(D)}\ 214 \qquad\textbf{(E)}\ 224$


 
 \textbf{Answer: }E



\textbf{Solution}

We have $161 = 7 \cdot 23.$ Since $a$ has two digits, the factors must be $23$ and $7,$ so $a = 32$ and $b = 7.$ Then, $ab = 7 \times 32 = \boxed{\mathrm{\textbf{(E)}\ } 224}.$

 



\textbf{Solution 2(cross-out/guess n check)}

We can cross out our answer choices. (B)161 would obviously not work, a quick look at (D)214 tells us that it only has two factor pairs which are both three digits, (A)116 is too low. This leaves (C)204 and (E)224. From here, we can simply guess and check to get (E)224.

 P.S. This solution is NOT advised, but is still listed as an alternate way.

 


	\item (\textit{AMC 10 10B 2011}) On Halloween Casper ate $1/3$ of his candies and then gave $2$ candies to his brother. The next day he ate $1/3$ of his remaining candies and then gave $4$ candies to his sister. On the third day he ate his final $8$ candies. How many candies did Casper have at the beginning?


 $\textbf{(A)}\ 30 \qquad\textbf{(B)}\ 39 \qquad\textbf{(C)}\ 48 \qquad\textbf{(D)}\ 57 \qquad\textbf{(E)}\ 66$


 
 \textbf{Answer: }A



\textbf{Solution 1}

Let $x$ represent the amount of candies Casper had at the beginning.

 
 \begin{align*} \frac{2}{3} \left(\frac{2}{3} x - 2\right) - 4 - 8 &= 0\\ \frac{2}{3} x - 2 &= 18\\ \frac{2}{3} x &= 20\\ x &= \boxed{\textbf{(A)} 30} \end{align*}


 



\textbf{Solution 2}

We work backwards. If he had 8 candies at the end, then before he gave candies to his sister he had 12 candies. This means that at the end of Halloween he had 18 candies, so before he gave candies to his brother he had 20 candies. Therefore, at the start he had $\boxed{\textbf{(A)} 30}$

 



\textbf{Solution 3 (answer choices)}

A solve by algebra is more secure and safe (and usually faster), but you can also test the answer choices. We are lucky and option A works $\Longrightarrow \boxed{\textbf{(A) } 30}$.

 


	\item (\textit{AMC 10 10B 2011}) The sum of two angles of a triangle is $6/5$ of a right angle, and one of these two angles is $30^{\circ}$ larger than the other. What is the degree measure of the largest angle in the triangle?


 $\textbf{(A)}\ 69 \qquad\textbf{(B)}\ 72 \qquad\textbf{(C)}\ 90 \qquad\textbf{(D)}\ 102 \qquad\textbf{(E)}\ 108$


 
 \textbf{Answer: }B



\textbf{Solution}

The sum of two angles in a triangle is $\frac{6}{5}$ of a right angle $\longrightarrow \frac{6}{5} \times 90 = 108$

 If $x$ is the measure of the first angle, then the measure of the second angle is $x+30$.
 \[x + x + 30 = 108 \longrightarrow 2x = 78 \longrightarrow x = 39\]

 Now we know the measure of two angles are $39^{\circ}$ and $69^{\circ}$. By the Triangle Sum Theorem, the sum of all angles in a triangle is $180^{\circ},$ so the final angle is $72^{\circ}$. Therefore, the largest angle in the triangle is $\boxed{\mathrm{(B) \ } 72}$

 


	\item (\textit{AMC 10 10B 2011}) At a certain beach if it is at least $80^{\circ} F$ and sunny, then the beach will be crowded. On June 10 the beach was not crowded. What can be concluded about the weather conditions on June 10?


 $\textbf{(A)}\ \text{The temperature was cooler than } 80^{\circ} \text{F and it was not sunny.}$


 $\textbf{(B)}\ \text{The temperature was cooler than } 80^{\circ} \text{F or it was not sunny.}$


 $\textbf{(C)}\ \text{If the temperature was at least } 80^{\circ} \text{F, then it was sunny.}$


 $\textbf{(D)}\ \text{If the temperature was cooler than } 80^{\circ} \text{F, then it was sunny.}$


 $\textbf{(E)}\ \text{If the temperature was cooler than } 80^{\circ} \text{F, then it was not sunny.}$


 
 \textbf{Answer: }B



\textbf{Solution}

The beach not being crowded only means that it is not both hot and sunny. Equivalently, it is either cool or cloudy, which means $\boxed{\textbf{(B)}}$ is correct.

 It could be in fact cool and sunny, which would violate \textbf{(A)} and \textbf{(E)}. It could also be cool and cloudy, which would violate \textbf{(D)}, or hot and cloudy, which would violate \textbf{(C)}.

 



\textbf{Solution 2 (Answer Choices)}

Looking at each of the answer choices:

 $\textbf{(A)}\ \text{The temperature was cooler than } 80^{\circ} \text{F and it was not sunny.}$ This doesn't include the fact that it can also be either cooler than $80^{\circ}$ and sunny or the temperature could be at least $80^{\circ}$ and not sunny, for the beach not to be crowded, so this is wrong.

 $\textbf{(B)}\ \text{The temperature was cooler than } 80^{\circ} \text{F or it was not sunny.}$ This is always correct, because for the beach to not be crowded,  at least 1 of the following conditions has to be true: 1. The temperature is cooler than $80^{\circ}$2. It is not sunny.Since $\textbf{(B)}$ has both of these, the answer is $\boxed{\textbf{(B)}}$.

 $\textbf{(C)}\ \text{If the temperature was at least } 80^{\circ} \text{F, then it was sunny.}$ This is obviously wrong, because if it was at least $80^{\circ}$, then it would have to be cloudy (or anything other than not-sunny), for there to not be a crowd.

 $\textbf{(D)}\ \text{If the temperature was cooler than } 80^{\circ} \text{F, then it was sunny.}$ This is wrong, because it forgets the fact that it can also be not-sunny, because the temperature is lower than $80^{\circ}$, so no matter whether it is sunny, or not-sunny, there will not be a crowd.

 $\textbf{(E)}\ \text{If the temperature was cooler than } 80^{\circ} \text{F, then it was not sunny.}$ The same said for $\textbf{(D)}$ goes for this choice too, because this statement forgets the fact that it can also be sunny, as it doesn't matter whether it is sunny or not, mainly because the temperature is not viable for the crowd to be in the beach.

 This concludes that the answer is $\boxed{\textbf{(B)}}$.

 


 


	\item (\textit{AMC 10 10B 2011}) 
The area of $\triangle$$EBD$ is one third of the area of $3-4-5$ $\triangle$$ABC$. Segment $DE$ is perpendicular to segment $AB$. What is $BD$? 

 


 
                    \begin{figure}[H]
                        \centering
                        \includegraphics[width=0.5\textwidth]{Problem/AMC_10_10B_2011_9_Problem_Figure_1.png}
                    \end{figure}
                    


 $\textbf{(A)}\ \frac{4}{3} \qquad\textbf{(B)}\ \sqrt{5} \qquad\textbf{(C)}\ \frac{9}{4} \qquad\textbf{(D)}\ \frac{4\sqrt{3}}{3} \qquad\textbf{(E)}\ \frac{5}{2}$


 
 \textbf{Answer: }D



\textbf{Solution 1}

$\triangle ABC \sim \triangle EBD$ by AA Similarity. Therefore $DE = \frac{3}{4} BD$. Find the areas of the triangles.
 \[\triangle ABC: 3 \times 4 \times \frac{1}{2} = 6\]
 \[\triangle EBD: BD \times \frac{3}{4} BD \times \frac{1}{2} = \frac{3}{8} BD ^2\]The area of $\triangle EBD$ is one third of the area of $\triangle ABC$.
 \begin{align*} \frac{3}{8} BD^2 &= 6 \times \frac{1}{3}\\ 9BD^2 &= 48\\ BD^2 &= \frac{16}{3}\\ BD &= \boxed{\textbf{(D)} \frac{4\sqrt{3}}{3}} \end{align*}


 



\textbf{Solution 2}

$\triangle ABC \sim \triangle EBD$ by AA Similarity. Since the area of $\triangle EBD$ is $\frac{1}{3}$ of $\triangle ABC$ and the bases/heights are in the same ratio, we use the formula forarea of a triangle for these ratios. Thus,
 \[\frac{1}{3}[\triangle ABC] = [\triangle EBD]\]
 \[\frac{1}{3} \times \frac{1}{2} \times AC \times BC = \frac{1}{2} \times ED \times DB\]In order to scale the sides of ED and DB to make $\frac{1}{3}$ (since the ratios of sides are the same), we take the square root of $\frac{1}{3} = \frac{\sqrt{3}}{3}$ to scale each side by the same amount. 

 Thus $BD = 4 \times \frac{\sqrt{3}}{3}$ and the answer is $BD = \boxed{\textbf{(D)} \frac{4\sqrt{3}}{3}}$

 



\textbf{Solution 3 (Shortcut)}

The ratio of the areas of $\triangle$$EBD$ and $\triangle$$ABC$ is $1 : 3$, meaning the ratio of the sides is $1 : \sqrt{3}$. The only answer choice involving $\sqrt{3}$ is $\boxed{\textbf{(D) } \frac{4\sqrt{3}}{3}}$.

 -Solution by Joeya

 


 



\textbf{Solution 4 (slightly more vigorous than Solution 3)}

The ratio of the areas is equal to twice the ratio of sides (in similar figures) because area is a second-degree property of similar figures. So like solution 3, the ratio of sides is $\sqrt{\frac{1}{3}}=\frac{1}{\sqrt{3}}=\sqrt{3}.$ By similarity, $\frac{ED}{DB}=\frac{\sqrt{3}}{DB}=\frac{3}{4}$, so solving for DB, we get $\boxed{\textbf{(D) } \frac{4\sqrt{3}}{3}}$.

 



\textbf{Solution 5}

It is apparent that $\Delta ABC \sim \Delta EBD$ by AA similarity ($\angle B=\angle B$ and $\angle EDB=\angle ACB$). Thus, let the side length of $ED$ equal $3x$ and $DB=4x.$ We can then see that $[EDB]=\dfrac{3x\cdot4x}2=6x^2$, and we are given that $[ABC]=3\cdot[EDB]$. Thus, $\dfrac{3\cdot4}2=3\cdot6x^2\implies6=18x^2\implies x=\dfrac{\sqrt{3}}3$. Since we let $BD=4x$, we know that $BD=\boxed{\textbf{(D) }\dfrac{4\sqrt{3}}3}$. ~Technodoggo

 


	\item (\textit{AMC 10 10B 2011}) Consider the set of numbers $\{1, 10, 10^2, 10^3, \ldots, 10^{10}\}$. The ratio of the largest element of the set to the sum of the other ten elements of the set is closest to which integer?


 $\textbf{(A)}\ 1 \qquad\textbf{(B)}\ 9 \qquad\textbf{(C)}\ 10 \qquad\textbf{(D)}\ 11 \qquad\textbf{(E)} 101$


 
 \textbf{Answer: }B



\textbf{Solution 1}

The requested ratio is 
 \[\dfrac{10^{10}}{10^9 + 10^8 + \ldots + 10 + 1}.\] Using the formula for a geometric series, we have 
 \[10^9 + 10^8 + \ldots + 10 + 1 = \dfrac{10^{10} - 1}{10 - 1} = \dfrac{10^{10} - 1}{9},\] which is very close to $\dfrac{10^{10}}{9},$ so the ratio is very close to $\boxed{\mathrm{(B) \ } 9}.$

 


 



\textbf{Solution 2}

The problem asks for the value of 
 \[\dfrac{10^{10}}{10^9 + 10^8 + \ldots + 10 + 1}.\] Written in base 10, we can find the value of $10^9 + 10^8 + \ldots + 10 + 1$ to be $1111111111.$ Long division gives us the answer to be $\boxed{\mathrm{(B) \ } 9}.$

 Alternate finish: multiply the denominator by 9 and notice that it is 1 less than $10^{10}$. So the answer is very very close to $\boxed{\mathrm{(B) } 9}$.

 



\textbf{Solution 3}

Let $f(n)=\dfrac{10^n}{1+10+10^2+10^3+\cdots+10^{n-1}}$. We are approximating $f(10)$. Trying several small values of $n$ gives answers very close to $9$, so our answer is $\boxed{\textbf{(B)}~9}$. Note that $f(1)=10$, but $f(2)=\dfrac{100}{11}\approx9.09.$ ~Technodoggo

 


	\item (\textit{AMC 10 10B 2011}) There are $52$ people in a room.  What is the largest value of $n$ such that the statement "At least $n$ people in this room have birthdays falling in the same month" is always true?


 $\textbf{(A)}\ 2 \qquad\textbf{(B)}\ 3 \qquad\textbf{(C)}\ 4 \qquad\textbf{(D)}\ 5 \qquad\textbf{(E)}\ 12$


 
 \textbf{Answer: }D



\textbf{Solution}

Pretend you have $52$ people you want to place  in $12$ boxes, because there are $12$ months in a year. By the \href{https://artofproblemsolving.com/wiki/index.php?title=Pigeonhole\_Principle}{Pigeonhole Principle}, one box must have at least $\left\lceil \frac{52}{12} \right\rceil$ people $\longrightarrow \boxed{\textbf{(D)} 5}$

 


	\item (\textit{AMC 10 10B 2011}) Keiko walks once around a track at exactly the same constant speed every day. The sides of the track are straight, and the ends are semicircles. The track has a width of $6$ meters, and it takes her $36$ seconds longer to walk around the outside edge of the track than around the inside edge. What is Keiko's speed in meters per second?


 


 

                    \begin{figure}[H]
                        \centering
                        \includegraphics[width=0.5\textwidth]{Problem/AMC_10_10B_2011_12_Problem_Figure_1.png}
                    \end{figure}
                    


 


 $\textbf{(A)}\ \frac{\pi}{3} \qquad\textbf{(B)}\ \frac{2\pi}{3} \qquad\textbf{(C)}\ \pi \qquad\textbf{(D)}\ \frac{4\pi}{3} \qquad\textbf{(E)}\ \frac{5\pi}{3}$


 
 \textbf{Answer: }A



\textbf{Solution 1}

Let $s$ be Keiko's speed in meters per second, $a$ be the length of the straight parts of the track, $b$ be the radius of the smaller circles, and $b+6$ be the radius of the larger circles. The length of the inner edge will be $2a+2b \pi$ and the length of the outer edge will be $2a+2\pi (b+6).$ Since it takes $36$ seconds longer for Keiko to walk on the outer edge,

 
 \begin{align*} \frac{2a+2b \pi}{s} + 36 &= \frac{2a+2\pi (b+6)}{s}\\ 2a+2b\pi +36s &= 2a+2b\pi +12\pi\\ 36s&=12\pi\\ s&=\boxed{\textbf{(A)} \frac{\pi}{3}} \end{align*}


 



\textbf{Solution 2}

It is basically the same as Solution 1 except you can completely disregard the straight edges of the track since it will take Keiko the same time to walk that length for both inner and outer circles. Instead, focus on the circular part. If the diameter of the smaller circle were $r,$ then the length of the smaller circle would be $r \pi$ and the length of the larger circle would be $(r+12) \pi.$ Since it still takes $36$ seconds longer,

 
 \begin{align*} \frac{r \pi}{s} + 36 &= \frac{(r+12)\pi}{s}\\ r \pi + 36s &= r \pi + 12 \pi\\ s&=\boxed{\textbf{(A)} \frac{\pi}{3}} \end{align*}


 


	\item (\textit{AMC 10 10B 2011}) Two real numbers are selected independently at random from the interval $[-20, 10]$. What is the probability that the product of those numbers is greater than zero?


 $\textbf{(A)}\ \frac{1}{9} \qquad\textbf{(B)}\ \frac{1}{3} \qquad\textbf{(C)}\ \frac{4}{9} \qquad\textbf{(D)}\ \frac{5}{9} \qquad\textbf{(E)}\ \frac{2}{3}$


 
 \textbf{Answer: }D



\textbf{Solution}

For the product of two numbers to be greater than zero, they either have to both be negative or both be positive. The interval for a positive number is $\frac{1}{3}$ of the total interval, and the interval for a negative number is $\frac{2}{3}$. Therefore, the probability the product is greater than zero is
 \[\frac{1}{3} \times \frac{1}{3} + \frac{2}{3} \times \frac{2}{3} = \frac{1}{9} + \frac{4}{9} = \boxed{\textbf{(D)} \frac{5}{9}}\]

 


 



\textbf{Solution 2}

We can plot this problem on the coordinate grid. First, plot the possible region [the square with endpoints $(10,10), (-20, 10), (-20, -20), (10, -20)$]. Including quadrant 3 and 4 allows for the second number selected to be negative. Then, we can find our possible region which is if both numbers are positive or if both numbers are negative. Lets take the first case in which both numbers are positive: The region which has these desired values is the square [$(10,10) ,(0,0) ,(10,0) ,(0,10)$]. Similarly if both numbers are negative we get the square [$(-20,-20), (0,0), (-20,0) ,(0,-20)$] Then we find the area of the possible region $30\cdot 30 = 900$ and then find the areas of our positive regions $10\cdot 10=100$  and $20\cdot 20 = 400$. Adding and simplifying we get $\frac{100+400}{900}=\frac{5}{9}$.

 



\textbf{Solution 3}

We will use complementary counting. The probability that the product is negative can be found by finding the probability that one number is positive and the other number is negative. The probability of a positive number being selected is $\frac13$, and the probability of a negative number being selected is $\frac23$. So, we see that the probability of the product being negative is
 \[2 \cdot \frac23 \cdot \frac13 = \frac49.\]

 We multiply by $2$ because you can either pick the negative or positive number first. 

 Thus, the probability of the product being positive is 
 \[1-\frac49 = \boxed{\textbf{(D) }\frac59}\]

 



\textbf{Solution 4}

We can simplify the problem down by dividing the boundary numbers by 10, which gives us $[-2,1]$. Then, it is easier for us to come to Solution 2. The probability the product is greater than zero is
 \[\frac{1}{3} \times \frac{1}{3} + \frac{2}{3} \times \frac{2}{3} = \frac{1}{9} + \frac{4}{9} = \boxed{\textbf{(D)} \frac{5}{9}}\]

 


	\item (\textit{AMC 10 10B 2011}) A rectangular parking lot has a diagonal of $25$ meters and an area of $168$ square meters. In meters, what is the perimeter of the parking lot?


 $\textbf{(A)}\ 52 \qquad\textbf{(B)}\ 58 \qquad\textbf{(C)}\ 62 \qquad\textbf{(D)}\ 68 \qquad\textbf{(E)}\ 70$


 
 \textbf{Answer: }C



\textbf{Solution}

Let the sides of the rectangular parking lot be $a$ and $b$. Then $a^2 + b^2 = 625$ and $ab = 168$. Add the two equations together, then factor.
 \begin{align*} a^2 + 2ab + b^2 &= 625 + 168 \times 2\\ (a + b)^2 &= 961\\ a + b &= 31 \end{align*}
The perimeter of a rectangle is $2 (a + b) = 2 (31) = \boxed{\textbf{(C)} 62}$

 



\textbf{Solution 2 (Quick)}

We see that the answer choices and the perimeter are integers. Therefore, the sides of the rectangle are most likely integers that satisfy $a^2+b^2=25^2$. In other words, $(a,b,25)$ is a set of Pythagorean triples. Guessing and checking, we have $(7,24,25)$ as the triplet, as the area is $7 \cdot 24 = 168$ as requested. Therefore, the perimeter is $2(7+24)=\boxed{\textbf{(C)} 62}$.

 


	\item (\textit{AMC 10 10B 2011}) Let $@$ denote the "averaged with" operation: $a @ b = (a+b)/2$. Which of the following distributive laws hold for all numbers $x, y,$ and $z$? 
 \[\text{I. x @ (y + z) = (x @ y) + (x @ z)}\] 
 \[\text{II. x + (y @ z) = (x + y) @ (x + z)}\] 
 \[\text{III. x @ (y @ z) = (x @ y) @ (x @ z)}\]


 $\textbf{(A)}\ \text{I only} \qquad\textbf{(B)}\ \text{II only} \qquad\textbf{(C)}\ \text{III only} \qquad\textbf{(D)}\ \text{I and III only} \qquad\textbf{(E)}\ \text{II and III only}$


 
 \textbf{Answer: }E



\textbf{Solution}

Simplify each operation and see which ones hold true.

 
 \begin{align*} \text{I.} \qquad x @ (y + z) &= (x @ y) + (x @ z)\\ \frac{x+y+z}{2} &= \frac{x+y}{2} + \frac{x+z}{2}\\ \frac{x+y+z}{2} &\not= \frac{2x+y+z}{2} \end{align*}


 
 \begin{align*} \text{II.} \qquad x + (y @ z) &= (x + y) @ (x + z)\\ x+ \frac{y+z}{2} &= \frac{2x+y+z}{2}\\ \frac{2x+y+z}{2} &= \frac{2x+y+z}{2} \end{align*}


 
 \begin{align*} \text{III.} \qquad x @ (y @ z) &= (x @ y) @ (x @ z)\\ x @ \frac{y+z}{2} &= \frac{x+y}{2} @ \frac{x+z}{2}\\ \frac{2x+y+z}{4} &= \frac{2x+y+z}{4} \end{align*}


 $\boxed{\textbf{(E)} \text{II and III only}}$

 



\textbf{Solution 2}

Alternatively, substitute arbitrary values for $x$, $y$, and $z$ where $x\neq y\neq z$. For example, $x=1$, $y=2$, and $z=3$ give that only II and III work. $\boxed{\textbf{(E)} \text{II and III only}}$

 


	\item (\textit{AMC 10 10B 2011}) A dart board is a regular octagon divided into regions as shown. Suppose that a dart thrown at the board is equally likely to land anywhere on the board. What is the probability that the dart lands within the center square?


 
                    \begin{figure}[H]
                        \centering
                        \includegraphics[width=0.5\textwidth]{Problem/AMC_10_10B_2011_16_Problem_Figure_1.png}
                    \end{figure}
                    


 $\textbf{(A)}\ \frac{\sqrt{2} - 1}{2} \qquad\textbf{(B)}\ \frac{1}{4} \qquad\textbf{(C)}\ \frac{2 - \sqrt{2}}{2} \qquad\textbf{(D)}\ \frac{\sqrt{2}}{4} \qquad\textbf{(E)}\ 2 - \sqrt{2}$


 
 \textbf{Answer: }A



\textbf{Solution}


                    \begin{figure}[H]
                        \centering
                        \includegraphics[width=0.5\textwidth]{Solution/AMC_10_10B_2011_16_Solution_Figure_1.png}
                    \end{figure}
                    

 If the side lengths of the dart board and the side lengths of the center square are all $\sqrt{2},$ then the side length of the legs of the triangles are $1$.

 
 \begin{align*} \text{area of center square} &: \sqrt{2} \times \sqrt{2} = 2\\ \text{total area} &: (\sqrt{2})^2 + 4(1 \times \sqrt{2}) + 4(1 \times 1 \times \frac{1}{2}) = 2 + 4\sqrt{2} + 2 = 4 + 4\sqrt{2} \end{align*}


 Use \href{https://artofproblemsolving.com/wiki/index.php?title=Geometric\_probability}{Geometric probability} by putting the area of the desired region over the area of the entire region.

 
 \[\frac{2}{4+4\sqrt{2}} = \frac{1}{2+2\sqrt{2}} \times \frac{2-2\sqrt{2}}{2-2\sqrt{2}} = \frac{2-2\sqrt{2}}{-4} = \boxed{\textbf{(A)} \frac{\sqrt{2}-1}{2}}\]

 
Explanation:The area of the octagon consists of the area of the triangles, the rectangles, and the square in the middle. Assume the octagon has side length $1$. The triangles are right isosceles triangles with the hypotenuse 1, so their side length is $\frac{\sqrt{2}}{2}$ and the area of one triangle is $\frac{\sqrt{2}}{2}\cdot\frac{\sqrt{2}}{2}\cdot\frac{1}{2}=\frac{1}{4}$. The area of all 4 triangles is then just $1$. The rectangles share one side with the octagon, and another with the triangle. The octagon has side length 1, and the triangle has side length $\frac{\sqrt{2}}{2}$ as found earlier. So the area of one rectangle is $\frac{\sqrt{2}}{2}$. The area of all 4 is $4\cdot\frac{\sqrt{2}}{2}=2\sqrt{2}$. Finally, the area of the square in the middle is $1\cdot1=1$. The total area is $1+1+2\sqrt{2}=2+2\sqrt{2}$. We want the area of the square over the area of the octagon, which is $\frac{1}{2+2\sqrt{2}}$. Rationalize by multiplying both numerator and denominator by $2-2\sqrt{2}$: $\frac{1}{2+2\sqrt{2}}\cdot\frac{2-2\sqrt{2}}{2-2\sqrt{2}}=\frac{2-2\sqrt{2}}{\left(2+2\sqrt{2}\right)\left(2-2\sqrt{2}\right)}$. By the difference of squares, the denominator reduces to $-4$ and the fraction is $\frac{2-2\sqrt{2}}{-4}=\frac{\sqrt{2}-1}{2}$ which is $\boxed{\textbf{(A) } \frac{\sqrt{2}-1}{2}}$.

 



\textbf{Solution 2}

The area of a regular octagon is $2(1+\sqrt{2})a^2$ where $a$ is the side. Hence the answer is obvious now.

 


	\item (\textit{AMC 10 10B 2011}) In the given circle, the diameter $\overline{EB}$ is parallel to $\overline{DC}$, and $\overline{AB}$ is parallel to $\overline{ED}$. The angles $AEB$ and $ABE$ are in the ratio $4 : 5$. What is the degree measure of angle $BCD$?


 
                    \begin{figure}[H]
                        \centering
                        \includegraphics[width=0.5\textwidth]{Problem/AMC_10_10B_2011_17_Problem_Figure_1.png}
                    \end{figure}
                    

 $\textbf{(A)}\ 120 \qquad\textbf{(B)}\ 125 \qquad\textbf{(C)}\ 130 \qquad\textbf{(D)}\ 135 \qquad\textbf{(E)}\ 140$


 
 \textbf{Answer: }C



\textbf{Solution 1}

We can let $\angle AEB$ be $4x$ and $\angle ABE$ be $5x$ because they are in the ratio $4 : 5$. When an \href{https://artofproblemsolving.com/wiki/index.php?title=Inscribed\_angle}{inscribed angle} contains the \href{https://artofproblemsolving.com/wiki/index.php?title=Diameter}{diameter}, the inscribed angle is a \href{https://artofproblemsolving.com/wiki/index.php?title=Right\_angle}{right angle}. Therefore by triangle sum theorem, $4x+5x+90=180 \longrightarrow x=10$ and $\angle ABE = 50$.

 $\angle ABE = \angle BED$ because they are alternate interior angles and $\overline{AB} \parallel \overline{ED}$. Opposite angles in a \href{https://artofproblemsolving.com/wiki/index.php?title=Cyclic}{cyclic} quadrilateral are \href{https://artofproblemsolving.com/wiki/index.php?title=Supplementary}{supplementary}, so $\angle BED + \angle BCD = 180$. Use substitution to get $\angle ABE + \angle BCD = 180 \longrightarrow 50 + \angle BCD = 180 \longrightarrow \angle BCD = \boxed{\textbf{(C)} 130}$

 \textbf{Note:} We could also tell that quadrilateral $BEDC$ is an isosceles trapezoid because for $\overline{EB}$ and $\overline{DC}$ to be parallel, the line going through the center of the circle and perpendicular to $\overline{DC}$ must fall through the center of $\overline{DC}$.

 



\textbf{Solution 2}

Note $\angle ABE = \angle BED=50$ as before. The sum of the interior angles for quadrilateral $EBCD$ is $360$. Denote the center of the circle as $P$. $\angle PDE = \angle PED = 50$. Denote $\angle PDC = \angle PCD = x$ and $\angle PBC = \angle PCB = y$. We wish to find $\angle BCD = x+y$. Our equation is $(\angle PDE +\angle PED)+(\angle PDC+\angle PCD)+(\angle PBC + \angle PCB) = 360 \longrightarrow 2(50) + 2x +2y = 360$. Our final equation becomes $2(x+y)+100 = 360$. After subtracting $100$ and dividing by $2$, our answer becomes $x+y=\boxed{\textbf{(C)} 130}$

 



\textbf{Solution 3 (Circle Geometry)}

Note that $\overset{\Large\frown} {BE}$ intercepts $\angle BAE$. Since, $\overset{\Large\frown} {BE}=180$, thus $\angle BAE=90°$ (courtesy of the Inscribed Angles Theorem). 

 Since we know that $\angle BAE=90°$, then $\angle AEB + \angle ABE = 90°$, (courtesy of the Triangle Sum Theorem) and also $5\angle AEB = 4\angle ABE$. By solving this variation, $\angle AEB = 40$ and $\angle ABE = 50$. After that, due to the Alternate Interior Angles Theorem, $\angle ABE \cong \angle BED$, which means $\angle BED = 50$. 

 
After doing some angle chasing, then these following facts should be true,$\overset{\Large\frown} {AB}=80$ $\overset{\Large\frown} {BD}=100$ $\overset{\Large\frown} {AE}=100$. 

 Note that the arcs have to equal 360, so, $360=\overset{\Large\frown} {AB}+\overset{\Large\frown} {BD}+\overset{\Large\frown} {DE}+\overset{\Large\frown} {AE}$

 $360=80+100+100+\overset{\Large\frown} {DE}$

 $\overset{\Large\frown} {DE}=80$

 Notice how $\overset{\Large\frown} {DB}$ intercepts $\angle BCD$ and that $\overset{\Large\frown} {DB}=\overset{\Large\frown} {DE}+\overset{\Large\frown} {AE}+\overset{\Large\frown} {AB}$. 

 $\overset{\Large\frown} {DB}=80+100+80$

 $\overset{\Large\frown} {DB}=260$

 According to the Inscribed Angles Theorem, $2\angle BCD=\overset{\Large\frown} {DB}$, therefore the answer is $\frac{260}{2}= \boxed{\textbf{(C)} 130}$

 


	\item (\textit{AMC 10 10B 2011}) Rectangle $ABCD$ has $AB = 6$ and $BC = 3$. Point $M$ is chosen on side $AB$ so that $\angle AMD = \angle CMD$. What is the degree measure of $\angle AMD$?


 $\textbf{(A)}\ 15 \qquad\textbf{(B)}\ 30 \qquad\textbf{(C)}\ 45 \qquad\textbf{(D)}\ 60 \qquad\textbf{(E)}\ 75$


 
 \textbf{Answer: }E



\textbf{Solution 1}


                    \begin{figure}[H]
                        \centering
                        \includegraphics[width=0.5\textwidth]{Solution/AMC_10_10B_2011_18_Solution_Figure_1.png}
                    \end{figure}
                    

 It is given that $\angle AMD \sim \angle CMD$. Since $\angle AMD$ and $\angle CDM$ are alternate interior angles and $\overline{AB} \parallel \overline{DC}$, $\angle AMD \cong \angle CDM \longrightarrow \angle CMD \cong \angle CDM$. Use the \href{https://artofproblemsolving.com/wiki/index.php?title=Base\_Angle\_Theorem}{Base Angle Theorem} to show $\overline{DC} \cong \overline{MC}$. We know that $ABCD$ is a \href{https://artofproblemsolving.com/wiki/index.php?title=Rectangle}{rectangle}, so it follows that $\overline{MC} = 6$. We notice that $\triangle BMC$ is a $30-60-90$ triangle, and $\angle BMC = 30^{\circ}$. If we let $x$ be the measure of $\angle AMD,$ then
 \begin{align*} 2x + 30 &= 180\\ 2x &= 150\\ x &= \boxed{\textbf{(E)} 75} \end{align*}


 


	\item (\textit{AMC 10 10B 2011}) What is the product of all the roots of the equation 
 \[\sqrt{5 | x | + 8} = \sqrt{x^2 - 16}.\]


 $\textbf{(A)}\ -64 \qquad\textbf{(B)}\ -24 \qquad\textbf{(C)}\ -9 \qquad\textbf{(D)}\ 24 \qquad\textbf{(E)}\ 576$


 
 \textbf{Answer: }A



\textbf{Solution 1}

First, square both sides, and isolate the absolute value.
 \begin{align*} 5|x|+8&=x^2-16\\ 5|x|&=x^2-24\\ |x|&=\frac{x^2-24}{5}. \\ \end{align*}
Solve for the absolute value and factor.

 Case 1: $x=\frac{x^2-24}{5}$

 Multiplying both sides by $5$ gives us 
 \[5x=x^2-24.\]Rearranging and factoring, we have
 \begin{align*} x^2-5x-24 &=0, \\ (x-8)(x+3) &= 0.\\ \end{align*}


 Case 2: $x=\frac{-x^2+24}{5}$

 As above, we multiply both sides by $5$ to find
 \[5x=-x^2+24.\]Rearranging and factoring gives us
 \begin{align*} x^2+5x-24 &=0, \\ (x+8)(x-3) &= 0. \\ \end{align*}


 Combining these cases, we have $x= -8, -3, 3, 8$. Because our first step of squaring is not reversible, however, we need to check for extraneous solutions. Plug each solution for $x$ back into the original equation to ensure it works. Whether the number is positive or negative does not matter since the absolute value or square will cancel it out anyways.Trying $|x|=|3|$, we have
 \begin{align*} \sqrt{5|3|+8}&=\sqrt{3^2-16}, \\ \sqrt{15+8}&=\sqrt{9-16}, \\ \sqrt{23} &\not= \sqrt{-7}.\\ \end{align*}
Therefore, $x = 3$ and $x= -3$ are extraneous.

 Checking $|x|=|8|$, we have
 \begin{align*} \sqrt{5|8|+8}&=\sqrt{8^2-16}, \\ \sqrt{40+8}&=\sqrt{64-16}, \\ \sqrt{48}&=\sqrt{48}.\\ \end{align*}


 The roots of our original equation are $-8$ and $8$ and product is $-8 \times 8 = \boxed{\textbf{(A)} -64}$.

 



\textbf{Solution 2}

Square both sides, to get $5|x| + 8 = x^2-16$. Rearrange to get $x^2 - 5|x| - 24 = 0$. Seeing that $x^2 = |x|^2$, substitute to get $|x|^2 - 5|x| - 24 = 0$. We see that this is a quadratic in $|x|$. Factoring, we get $(|x|-8)(|x|+3) = 0$, so $|x| = \{8,-3\}$. Since the radicand of the equation can't be negative, the sole solution is $|x| = 8$. Therefore, $x$ can be $8$ or $-8$. The product is then $-8 \times 8 = \boxed{\textbf{(A)} -64}$.

 



\textbf{Solution 3}

First we note that $x \in (-\infty,-4] \cup [4,\infty]$. This will help us later with finding extraneous solutions. Next, we have two cases:
 \[\text{IF } x\leq 4:\]
 \[\text{ } \sqrt{-5x+8} = \sqrt{x^2-16} \implies x^2+5x-24=0 \rightarrow x = -8,3\].We note that $3$ is not in the range of possible $x$'s and thus is not a solution.

 
 \[\text{IF } x\geq 4:\]
 \[\text{ } \sqrt{5x+8} = \sqrt{x^2-16} \implies x^2-5x-24=0 \rightarrow x = -3,8\].We again not that $-3$ is an extraneous solution.

 Thus, we have the two solutions $-8$ and $8$. Therefore product is $-8 \cdot 8 = \boxed{\textbf{(A)} -64}$.

 -ConfidentKoala4

 



\textbf{Solution 4}

To make this problem easier to comprehend, we can define variable $a = |x|$, with the condition that $a$ is always nonnegative. Also, since any number squared is always nonnegative, we can define $|x|^2 = x^2$. Then we can square both sides and substitute:
 \[\sqrt{5 | x | + 8} = \sqrt{x^2 - 16} => 5a + 8 = a^2 - 16.\]Bringing the equation over to one side, we get:
 \[a^2 - 5a - 24 = 0.\]Solving for a by factoring, we get:
 \[a^2 - 5a - 24 = 0 = (a - 8)(a + 3) = 0,\]so $a = 8$ or $a = -3$. But $a$ must be nonnegative, so the only value that works is $8$. If $a = 8$, then $|x| = 8$ and $x^2 = 8^2 = 64$, so $x$ can equal $8 or -8$. Multiplying the two, we get $8 * -8 = \boxed{\textbf{(A)} -64}$.

 


	\item (\textit{AMC 10 10B 2011}) Rhombus $ABCD$ has side length $2$ and $\angle B = 120^\circ$. Region $R$ consists of all points inside the rhombus that are closer to vertex $B$ than any of the other three vertices. What is the area of $R$?


 $\textbf{(A)}\ \frac{\sqrt{3}}{3} \qquad\textbf{(B)}\ \frac{\sqrt{3}}{2} \qquad\textbf{(C)}\ \frac{2\sqrt{3}}{3} \qquad\textbf{(D)}\ 1 + \frac{\sqrt{3}}{3} \qquad\textbf{(E)}\ 2$


 
 \textbf{Answer: }C



\textbf{Solution 1}

Suppose that $P$ is a point in the rhombus $ABCD$ and let $\ell_{BC}$ be the \href{https://artofproblemsolving.com/wiki/index.php?title=Perpendicular\_bisector}{perpendicular bisector} of $\overline{BC}$. Then $PB < PC$ if and only if $P$ is on the same side of $\ell_{BC}$ as $B$. The line $\ell_{BC}$ divides the plane into two half-planes; let $S_{BC}$ be the half-plane containing $B$. Let us define similarly $\ell_{BD},S_{BD}$ and $\ell_{BA},S_{BA}$. Then $R$ is equal to $ABCD \cap S_{BC} \cap S_{BD} \cap S_{BA}$. The region turns out to be an irregular pentagon. We can make it easier to find the area of this region by dividing it into four triangles:

 
                    \begin{figure}[H]
                        \centering
                        \includegraphics[width=0.5\textwidth]{Solution/AMC_10_10B_2011_20_Solution_Figure_1.png}
                    \end{figure}
                    Since $\triangle BCD$ and $\triangle BAD$ are equilateral, $\ell_{BC}$ contains $D$, $\ell_{BD}$ contains $A$ and $C$, and $\ell_{BA}$ contains $D$. Then $\triangle BEF \cong \triangle BGF \cong \triangle BGH \cong \triangle BIH$ with $BE = 1$ and $EF = \frac{1}{\sqrt{3}}$ so $[BEF] = \frac{1}{2}\cdot 1 \cdot \frac{\sqrt{3}}{3}$. Multiply this by 4 and it turns out that the pentagon has area $\boxed{(C)\frac{2\sqrt{3}}{3}}$.

 



\textbf{Solution 2}

We follow the steps shown above until we draw pentagon $BIHFE$. We know that rhombus $ABCD$ can be divided into equilateral triangles $\triangle ABD$ and $\triangle CBD$. Using the $30-60-90$ special right triangle rules, we find the height of the equilateral triangles (and the height of the rhombus) to be $\sqrt{3}$. Therefore, the area of $ABCD$ is $2\sqrt{3}$. We now have to take off the areas $\triangle CDA$, $\triangle CEF$, and $\triangle AIH$ to get the desired shape. $\triangle CDA$ is just half of $ABCD$ $(\sqrt {3})$ and $\triangle AIH$ and $\triangle CEF$ are each $\frac{\sqrt {3}}{6}$, for a total area of $2\sqrt {3}-\sqrt {3}-\frac{\sqrt{3}}{6}-\frac{\sqrt{3}}{6}=\boxed{(C)\frac{2\sqrt{3}}{3}}$.

 



\textbf{Solution 3}

We split rhombus $ABCD$ into two equilateral triangles, $ABD$ and $BCD$. In triangle $ABD$, the probability that a randomly selected point is closer to $B$ than either other point is $\frac{1}{3}$ *. Similarly, in triangle $BCD$, the same principle applies. Thus, the area of the region closer to $B$ than $A$, $C$, or $D$ is $\frac{1}{3} [ABD] + \frac{1}{3} [BCD]$. Since $ABD$ and $BCD$ are congruent, we have $\frac{1}{3} [ABD] + \frac{1}{3} [BCD] = \frac{2}{3} [ABD] = \frac{2}{3} \cdot \frac{s^2\sqrt{3}}{4} = \frac{2}{3} \cdot \frac{(2)^2\sqrt{3}}{4} = \boxed{\frac{2\sqrt3}{3} = C}$, and we are done.

 \*See note

 



\textbf{Solution 4}


                    \begin{figure}[H]
                        \centering
                        \includegraphics[width=0.5\textwidth]{Solution/AMC_10_10B_2011_20_Solution_Figure_2.png}
                    \end{figure}
                    Since $H$ and $E$ are halfway between $AB$ and $CB$, respectively, we know that $\overline{BH}=\overline{BE}=1$. By symmetry, $\Delta BFG$ is equilateral, so $\angle FBG=60^\circ\implies\angle EBF=\angle HBG=30^\circ$ and therefore $\Delta EBF$ and $\Delta HBG$ are 30-60-90 right triangles. Thus, $[\Delta EBF]=[\Delta BFG]=\dfrac1{2\sqrt3}$. We know that $\overline{FB}=\overline{GB}=\dfrac2{\sqrt3}$, so therefore $[\Delta BFG]=\dfrac{\sqrt3}4\left(\dfrac2{\sqrt3}\right)^2=\dfrac1{\sqrt3}$. Summing these three regions, we get $\dfrac1{2\sqrt3}+\dfrac1{2\sqrt3}+\dfrac1{\sqrt3}=\boxed{\textbf{(C)}~\dfrac{2\sqrt3}3}$. ~ Technodoggo, Asymptote diagram modified from Solution 1

 



\textbf{Solution 5}


                    \begin{figure}[H]
                        \centering
                        \includegraphics[width=0.5\textwidth]{Solution/AMC_10_10B_2011_20_Solution_Figure_3.png}
                    \end{figure}
                    To keep it simple, break rhombus $ABCD$ into two triangles, $ABD$ and $BCD$. To see the area closest to the point $B$, notice that a third of each triangle, which contains all the points nearest to $B$ in each triangle, is easily visualizable. Thus, a third of rhombus $ABCD$ must be found.

 We find the total area of rhombus $ABCD$, which we can again split into two congruent equilateral triangles with side length $2$. Using the formula of equilateral triangles and then multiplying by $\dfrac{1}{3}$:
 \[\dfrac{\sqrt{3}}{4}\cdot 2^2 \cdot 2 \cdot \dfrac{1}{3} = 2\sqrt{3} \cdot \dfrac{1}{3} = \boxed{\textbf{(C)}~\dfrac{2\sqrt{3}}{3}}\]-NSAoPS, diagram modified from Solution 1.

 \textbf{Note} There is an intuitive explanation for why the shaded region is one third of the rhombus. Consider one of the equilateral triangles; the figure is symmetrical, so one half of the shaded region should be in the triangle. Note that the perpendicular bisectors we constructed in previous solutions to deliminate $R$ must be medians, as the triangle is equilateral. Therefore, recalling that the three intersecting medians of a triangle divide it into 6 smaller triangles of equal area, we observe that the shaded region in the equilateral triangles includes precisely two of these smaller triangles. We can repeat this with the other equilateral triangle that forms the rhombus, and thus we conclude that the shaded area is one third of the rhombus's total area. 

 


	\item (\textit{AMC 10 10B 2011}) Brian writes down four integers $w > x > y > z$ whose sum is $44$. The pairwise positive differences of these numbers are $1, 3, 4, 5, 6,$ and $9$. What is the sum of the possible values for $w$?


 $\textbf{(A)}\ 16 \qquad\textbf{(B)}\ 31 \qquad\textbf{(C)}\ 48 \qquad\textbf{(D)}\ 62 \qquad\textbf{(E)}\ 93$


 
 \textbf{Answer: }B



\textbf{Solution 1}

The largest difference, $9,$ must be between $w$ and $z.$

 The smallest difference, $1,$ must be directly between two integers. This also means the differences directly between the other two should add up to $8.$ The only remaining differences that would make this possible are $3$ and $5.$ However, those two differences can't be right next to each other because they would make a difference of $8,$ which isn't given as a possibility in the problem. This means $1$ must be the difference between $y$ and $x.$ We can express the possible configurations as the lines.

 


 
                    \begin{figure}[H]
                        \centering
                        \includegraphics[width=0.5\textwidth]{Solution/AMC_10_10B_2011_21_Solution_Figure_1.png}
                    \end{figure}
                    

 If we look at the first number line, you can express $x$ as $w-5,$ $y$ as $w-6,$ and $z$ as $w-9.$ Since the sum of all these integers equal $44$,
 \begin{align*} w+w-5+w-6+w-9&=44\\ 4w&=64\\ w&=16 \end{align*}
You can do something similar to this with the second number line to find the other possible value of $w.$
 \begin{align*} w+w-3+w-4+w-9&=44\\ 4w&=60\\ w&=15 \end{align*}
The sum of the possible values of $w$ is $16+15 = \boxed{\textbf{(B) }31}$

 



\textbf{Solution 2}

First, like Solution 1, we know that $w-z=9 \ \text{(1)}$, because no two numbers could have a larger difference. Next, we find the sum of all the differences; since $w$ is in the positive part of a difference 3 times, and has no differences where it contributes as the negative part, the sum of the differences includes $3w$. Continuing in this way, we find that 
 \[3w+x-y-3z=28 \ \text{(2)}\]. Now, we can subtract $3w-3z=27$ from (2) to get $x-y=1 \ \text{(3)}$. Also, adding (2) with $w+x+y+z=44$ gives $4w+2x-2z=72$, or $2w+x-z=36$. Subtracting (1) from this gives $w+x=27$. Since we know $w-z$ and $x-y$, we find that 
 \[(w-z)+(x-y)=(w-y)+(x-z)=9+1=10\]. This means that $w-y$ and $x-z$ must be 4 and 6, in some order. If $w-y=6$, then subtracting this from (3) gives $(w-y)-(x-y)=6-1=5$, so $w-x=5$. This means that $(w-x)+(w+x)=2w=27+5=32$, so $w=16$. Similarly, $w$ can also equal $15$.

 Now if you are in a rush, you most likely would have answered $16+15=\boxed{\textbf{(B) }31}$. But we do have to check if these work. In fact, they do, giving solutions $(w,x,y,z)=(16, 11, 10, 7)$ and $(w,x,y,z)=(15, 12, 11, 6)$.

 



\textbf{Solution 3}

Let $w - x = a$, $w - y = b$, $w - z = c$. As above, we know that $c = 9$. Thus, $a < b < c$. So, we have $w + x + y + z = w + (w - a) + (w - b) + (w - 9) = 4w - a - b - 9 = 44$. This means $a + b + 9$ is a multiple of $4$. Testing values of $a$ and $b$, we find $(a, b, c) = (1, 6, 9), (3, 4, 9),$ and $(5, 6, 9)$ all satisfy this relation. The corresponding $(w, x, y, z)$ sets are $(15, 14, 9, 6), (15, 12, 11, 6),$ and $(16, 11, 10, 7)$. The first set does not satisfy the given conditions, but the other two do. Thus, $w = 15$ and $w = 16$ are both possible solutions so the answer is $16+15=\boxed{\textbf{(B) }31}$.

 



\textbf{Solution 4}

From the problem, we know that $w+x+y+z=44$. Since it is said that the pairwise positive differences between numbers are 1, 3, 4, 5, 6, 9, and we can figure that the pairwise positive differences are $(w-x)$, $(w-y)$, $(w-z)$, $(x-y)$, $(x-z)$, $(y-z)$, the sum of $(w-x), (w-y), (w-z), (x-y), (x-z), (y-z)$ is equal to the sum of 1, 3, 4, 5, 6, 9, so $(w-x)+(w-y)+(w-z)+(x-y)+(x-z)+(y-z)=28$. Simplifying, we get $3w-3z+x-y=28$. Adding $w+x+y+z=44$ and $(w-x)+(w-y)+(w-z)+(x-y)+(x-z)+(y-z)=28$, we get $4w+2x-2z=72$, and simplifying we get $2w+x-z=36$. Since $x-z$ is one of our positive differences, we can start guessing values for $w$, and if the equation simplifies to one of our numerical positive differences, that value of $w$ should work. We can start at $w=18$ and keep going down, because our sum has to be positive. For $w=18$, $x-z=0$, which is not one of our sums. For $w=17$, $x-z=2$,which is not one of our sums. For $w=16$, $x-z=4$, which is one of our sums, so 16 works. For $w=15$, $x-z=6$, which is one of our sums, so 15 works. For $w=14$, $x-z=8$, which is not one of our sums. If we keep going, $x-z$ will soon exceed 10 and exceed all our sums, so any value below $w=14$ will not work. Therefore, our only solutions for $w$ are 15 and 16, which means our sum is $\boxed{\textbf{(B) }31}$. You can check that 15 and 16 work by forming a string of 4 numbers as shown above.

 



\textbf{Solution 5}

Because we know that $w>x>y>z$ and that the positive differences are $1, 3, 4, 5, 6, 9$, we can immediately come to the conclusion that $w-z= 9$ (because w is the largest integer and z is the smallest integer, so their difference must be the greatest). With this we have $3$ equations, we have that $x+y+z+w=44$ (from the problem), $w-z=9$ and because we can add up all the possible possible differences (as shown in the previous solutions), we get that $3w+x-y-3z=28$. With these equations, we eventually manipulate these equations by doing the first equation minus 3 times the second to get $x-y=1$. We can also add the first and third equation and subtract the second equation to get $w+x=27$ thus we know that $w-x$ can be $3, 4, 5,$ and $6$ (note: 1 and 9 are not possible because we know that $x-y=1$ and $w-z=9$ and the remaining differences can only be taken by 1 pair, so $w-x$ cannot be equal to 1 or 9). However, in order to get an integer value for x and z, we find that $w-z$ can only be equal to 3 and 5. Thus, by solving these, we see that $w= 15$ and $w=16$. $15+16=\boxed{\textbf{(B) }31}$.

 


	\item (\textit{AMC 10 10B 2011}) A pyramid has a square base with sides of length $1$ and has lateral faces that are equilateral triangles. A cube is placed within the pyramid so that one face is on the base of the pyramid and its opposite face has all its edges on the lateral faces of the pyramid. What is the volume of this cube?


 $\textbf{(A)}\ 5\sqrt{2} - 7 \qquad\textbf{(B)}\ 7 - 4\sqrt{3} \qquad\textbf{(C)}\ \frac{2\sqrt{2}}{27} \qquad\textbf{(D)}\ \frac{\sqrt{2}}{9} \qquad\textbf{(E)}\ \frac{\sqrt{3}}{9}$


 
 \textbf{Answer: }A



\textbf{Solution 1}

It is often easier to first draw a diagram for such a problem.

 
                    \begin{figure}[H]
                        \centering
                        \includegraphics[width=0.5\textwidth]{Solution/AMC_10_10B_2011_22_Solution_Figure_1.png}
                    \end{figure}
                    Sometimes, it may also be easier to think of the problem in 2D. Take a cross section of the pyramid through the apex and two points from the base that are opposite to each other. Place it in two dimensions.

 
                    \begin{figure}[H]
                        \centering
                        \includegraphics[width=0.5\textwidth]{Solution/AMC_10_10B_2011_22_Solution_Figure_2.png}
                    \end{figure}
                    

 The dimensions of this triangle are $1, 1,$ and $\sqrt{2}$ because the sidelengths of the pyramid are $1$ and the base of the triangle is the diagonal of the pyramid's base. This is a $45-45-90$ triangle. Also, we can let the dimensions of the rectangle be $s$ and $s\sqrt{2}$ because the longer length was the diagonal of the cube's base and the shorter length was a side of the cube.

 
                    \begin{figure}[H]
                        \centering
                        \includegraphics[width=0.5\textwidth]{Solution/AMC_10_10B_2011_22_Solution_Figure_3.png}
                    \end{figure}
                    

 The two triangles on the right and left of the rectangle are also $45-45-90$ triangles because the rectangle is perpendicular to the base, and they share a $45^\circ$ angle with the larger triangle. Therefore, the legs of the right triangles can be expressed as $s.$

 
                    \begin{figure}[H]
                        \centering
                        \includegraphics[width=0.5\textwidth]{Solution/AMC_10_10B_2011_22_Solution_Figure_4.png}
                    \end{figure}
                    

 Now we can just use segment addition to find the value of $s.$
 \[\sqrt{2}=s+s\sqrt{2}+s=2s+s\sqrt{2}=(2+\sqrt{2})s\]
 \[s=\frac{\sqrt{2}}{2+\sqrt{2}}=\frac{1}{\sqrt{2}+1}=\frac{\sqrt{2}-1}{2-1}=\sqrt{2}-1\]

 The volume of the cube is $s^3 = (\sqrt{2}-1)^3 = (\sqrt{2}-1)(3-2\sqrt{2}) = 3\sqrt{2}-3-4+2\sqrt{2} = \boxed{\textbf{(A)} 5\sqrt{2}-7}$

 



\textbf{Solution 2}

Let $s$ be the slant height and $d$ be the distance from one side of the base to the point right under the apex. Then using the Pythagorean Theorem, the altitude from the apex to the base is 
 \[a=\sqrt{s^2-d^2}=\sqrt{\left(\dfrac{\sqrt{3}}{2}\right)^2-\left(\dfrac{1}{2}\right)^2}=\dfrac{\sqrt{2}}{2}\]Notice that the smaller pyramid on top of the cube is similar to the larger pyramid. Thus, letting $x$ be the edge length of the cube, then the altitude of the smaller pyramid is $\dfrac{\sqrt{2}}{2}\cdot{\dfrac{x}{1}}=\dfrac{\sqrt{2}}{2}x$.

 Since the altitude of the larger pyramid is the sum of the edge length of the cube and the altitude of the smaller pyramid, we have 
 \[\dfrac{\sqrt{2}}{2}=x+\dfrac{\sqrt{2}}{2}x\] 
 \[x=\sqrt{2}-1\]

 Finally, $V_{cube}=x^3=(\sqrt{2}-1)^3=\boxed{\textbf{(A)} 5\sqrt{2}-7}$

 



\textbf{Solution 3(basics)}

The lateral faces are equilateral triangles with side length $1$, so the slant height is $\frac{\sqrt{3}}{2}$. The distance from the base center to a vertex is $\frac{\sqrt{2}}{2}$, so the pyramid's altitude is

 
 \[h = \sqrt{1^2 - \left(\frac{\sqrt{2}}{2}\right)^2} = \frac{\sqrt{2}}{2}.\]

 
In the diagonal cross-section (base $\sqrt{2}$, height $\frac{\sqrt{2}}{2}$), let $x$ be half the cube’s side length. The alignment condition gives

 
 \[(0.5 - x)\sqrt{2} = 2x.\]

 Solving:

 
 \[0.5\sqrt{2} = x(2 + \sqrt{2}) \quad \Rightarrow \quad x = \frac{\sqrt{2} - 1}{2}.\]

 The cube’s side length is $2x = \sqrt{2} - 1$, so the volume is

 
 \[(2x)^3 = (\sqrt{2} - 1)^3 = 5\sqrt{2} - 7\]

 
 \[\boxed{5\sqrt{2} - 7}\]

 by Iamokatmath

 


	\item (\textit{AMC 10 10B 2011}) What is the hundreds digit of $2011^{2011}$?


 $\textbf{(A)}\ 1 \qquad\textbf{(B)}\ 4 \qquad\textbf{(C)}\ 5 \qquad\textbf{(D)}\ 6 \qquad\textbf{(E)}\ 9$


 
 \textbf{Answer: }D



\textbf{Solution 1}

Since $2011 \equiv 11  \pmod{1000},$ we know that $2011^{2011} \equiv 11^{2011}  \pmod{1000}.$

 To compute this, we use a clever application of the \href{https://artofproblemsolving.com/wiki/index.php?title=Binomial\_theorem}{binomial theorem}.

 $\begin{aligned} 11^{2011} &= (1+10)^{2011} \\ &= 1 + \dbinom{2011}{1} \cdot 10 + \dbinom{2011}{2} \cdot 10^2 + \cdots \end{aligned}$

 In all of the other terms, the power of $10$ is greater than $2$ and so is equivalent to $0$ modulo $1000,$ which means we can ignore it. We have:

 $\begin{aligned}11^{2011} &\equiv 1 + 2011\cdot 10 + \dfrac{2011 \cdot 2010}{2} \cdot 100 \\ &\equiv 1+20110 + \dfrac{11\cdot 10}{2} \cdot 100\\ &= 1 + 20110 + 5500\\ &\equiv 1 + 110 + 500\\&=611 \pmod{1000} \end{aligned}$

 Therefore, the hundreds digit is $\boxed{\textbf{(D) } 6}.$

 Side note: By \href{https://artofproblemsolving.com/wiki/index.php?title=Euler\%27s\_Totient\_Theorem}{Euler's Totient Theorem}, $a^{\phi (1000)} \equiv 1 \pmod{1000}$ 

 


	\item (\textit{AMC 10 10B 2011}) A lattice point in an $xy$-coordinate system is any point $(x, y)$ where both $x$ and $y$ are integers. The graph of $y = mx +2$ passes through no lattice point with $0 < x \le 100$ for all $m$ such that $1/2 < m < a$. What is the maximum possible value of $a$?


 $\textbf{(A)}\ \frac{51}{101} \qquad\textbf{(B)}\ \frac{50}{99} \qquad\textbf{(C)}\ \frac{51}{100} \qquad\textbf{(D)}\ \frac{52}{101} \qquad\textbf{(E)}\ \frac{13}{25}$


 
 \textbf{Answer: }B



\textbf{Solution 1}

For $y=mx+2$ to not pass through any lattice points with $0<x\leq 100$ is the same as saying that $mx\notin\mathbb Z$ for $x\in\{1,2,\dots,100\}$, or in other words, $m$ is not expressible as a ratio of positive integers $s/t$ with $t\leq 100$. Hence the maximum possible value of $a$ is the first real number after $1/2$ that is so expressible.

 For each $d=2,\dots,100$, the smallest multiple of $1/d$ which exceeds $1/2$ is $1,\frac23,\frac34,\frac35,\dots,\frac{50}{98},\frac{50}{99},\frac{51}{100}$ respectively, and the smallest of these is $\boxed{\textbf{(B)}\frac{50}{99}}$.

 



\textbf{Solution 1.1}

Note that finding the next possible slope $m$ greater than $\frac{1}{2}$ is the same as finding the next Farey fraction of order $100$ after $\frac{1}{2}$. Hence, we can use the Farey mediant theorem (\href{https://en.wikipedia.org/wiki/Farey\_sequence\#Farey\_neighbours}{https://en.wikipedia.org/wiki/Farey\_sequence\#Farey\_neighbours}), which states that if $\frac{a}{b}<\frac{c}{d}$ are neighboring Farey fractions of order $n$, then $bc-ad=1$ where $b, d \leq n$.

 In this case, we have $a=1$, $b=2$, and so $2c-d=1$. We try $d=100$, but that doesn't give integer $c$. We try $d=99$, which gives us $c=50$ so the upper bound is $\boxed{\textbf{(B)}\frac{50}{99}}$.

 



\textbf{Solution 2}

We see that for the graph of $y=mx+2$ to not pass through any lattice points, the denominator of $m$ must be greater than $100$, or else it would be canceled by some $0<x\le100$ which would make $y$ an integer. By using common denominators, we find that the order of the fractions from smallest to largest is $\text{(A), (B), (C), (D), (E)}$. We can see that when $m=\frac{50}{99}$, $y$ could be an integer, so therefore any fraction greater than $\frac{50}{99}$ would not work, as substituting our fraction $\frac{50}{99}$ for $m$ would produce an integer for $y$. So now we are left with only $\frac{51}{101}$ and $\frac{50}{99}$. But since $\frac{51}{101}=\frac{5049}{9999}$ and $\frac{50}{99}=\frac{5050}{9999}$, we can be absolutely certain that there isn't a number between $\frac{51}{101}$ and $\frac{50}{99}$ that can reduce to a fraction whose denominator is less than or equal to $100$. Since we are looking for the maximum value of $a$, we take the larger of $\frac{51}{101}$ and $\frac{50}{99}$, which is $\boxed{\textbf{(B)}\frac{50}{99}}$.

 



\textbf{Solution 3}

We want to find the smallest $m$ such that there will be an integral solution to $y=mx+2$ with $0<x\le100$. We first test A, but since the denominator has a $101$, $x$ must be a nonzero multiple of $101$, but it then will be greater than $100$. We then test B. $y=\frac{50}{99}x+2$ yields the solution $(99,52)$ which satisfies $0<x\le100$. Checking the answer choices, we know that the largest possible $a$ must be $\frac{50}{99}\implies\boxed{\textbf{(B)}}$

 



\textbf{Solution 4}

Notice that for $y=\frac{1}{2}x+2=\frac{50}{100}x+2$, $x=99$ is one of the integral values of $x$ such that the value of $\frac{50}{100}x$ is the closest to its next integral value.

 
Thus the maximum value for $a$ is the value of $m$ when the equation $y=99m+2$ goes through its next lattice point, which occurs when $m=\frac{b}{99}$ for some positive integer $b$.

 
Finding the common denominator, we have 
 \[\frac{50}{100}=\frac{4950}{9900}, \frac{b}{99}=\frac{100b}{9900}\] Since $a>\frac{1}{2}$, the smallest value for $b$ such that $100b>4950$ is $b=50$. 

 Thus the maximum value of $a$ is $\frac{50}{99}.\boxed{\mathrm{(B)}}$

 


 



\textbf{Solution 5}

For $0<x\leq 100$, the nearest lattice point directly above the line $y=\frac{1}{2}x+2$ is $(x,\frac{1}{2}x+3)$ if $x$ is even and $(x,\frac{1}{2}x+\frac{5}{2})$ if $x$ is odd. The slope of the line that contains this point and $(0,2)$ is $\frac{1}{2}+\frac{1}{x}$ if $x$ is even and $\frac{1}{2}+\frac{1}{2x}$ if $x$ is odd. The minimum value of the slope is $\frac{51}{100}$ if $x$ is even and $\frac{50}{99}$ if $x$ is odd. Therefore the line $y=mx+2$ contains no lattice point with $0<x\leq 100$ for $\frac{1}{2}<m<\boxed{\textbf{(B)}\frac{50}{99}}$.

 


	\item (\textit{AMC 10 10B 2011}) Let $T_1$ be a triangle with sides $2011, 2012,$ and $2013$. For $n \ge 1$, if $T_n = \triangle ABC$ and $D, E,$ and $F$ are the points of tangency of the incircle of $\triangle ABC$ to the sides $AB, BC$ and $AC,$ respectively, then $T_{n+1}$ is a triangle with side lengths $AD, BE,$ and $CF,$ if it exists. What is the perimeter of the last triangle in the sequence $( T_n )$?


 $\textbf{(A)}\ \frac{1509}{8} \qquad\textbf{(B)}\ \frac{1509}{32} \qquad\textbf{(C)}\ \frac{1509}{64} \qquad\textbf{(D)}\ \frac{1509}{128} \qquad\textbf{(E)}\ \frac{1509}{256}$


 
 \textbf{Answer: }D



\textbf{Solution 1}

By constructing the bisectors of each angle and the perpendicular radii of the incircle the triangle consists of 3 kites. 

 
                    \begin{figure}[H]
                        \centering
                        \includegraphics[width=0.5\textwidth]{Solution/AMC_10_10B_2011_25_Solution_Figure_1.png}
                    \end{figure}
                    Hence $AD=AF$ and $BD=BE$ and $CE=CF$. Let $AD = x, BD = y$ and $CE = z$ gives three equations:

 $x+y = a-1$

 $x+z = a$

 $y+z = a+1$

 (where $a = 2012$ for the first triangle.)

 Solving gives:

 $x= \frac{a}{2} -1$

 $y = \frac{a}{2}$

 $z = \frac{a}{2}+1$

 Subbing in gives that $T_2$ has sides of $1005, 1006, 1007$.

 $T_3$ can easily be derived from this as the sides still differ by 1 hence the above solutions still work (now with $a=1006$). All additional triangles will differ by one as the solutions above differ by one so this process can be repeated indefinitely until the side lengths no longer form a triangle.

 Subbing in gives $T_3$ with sides $502, 503, 504$.

 $T_4$ has sides $\frac{501}{2}, \frac{503}{2}, \frac{505}{2}$.

 $T_5$ has sides $\frac{499}{4}, \frac{503}{4}, \frac{507}{4}$.

 $T_6$ has sides $\frac{495}{8}, \frac{503}{8}, \frac{511}{8}$.

 $T_7$ has sides $\frac{487}{16}, \frac{503}{16}, \frac{519}{16}$.

 $T_8$ has sides $\frac{471}{32}, \frac{503}{32}, \frac{535}{32}$.

 $T_9$ has sides $\frac{439}{64}, \frac{503}{64}, \frac{567}{64}$.

 $T_{10}$ has sides $\frac{375}{128}, \frac{503}{128}, \frac{631}{128}$.

 $T_{11}$ would have sides $\frac{247}{256}, \frac{503}{256}, \frac{759}{256}$ but these lengths do not make a triangle as 
 \[\frac{247}{256} + \frac{503}{256} < \frac{759}{256}\]

 
Likewise, you could create an equation instead of listing all the triangles to $T_{11}$.The sides of a triangle $T_{k}$ would be 
 \[\frac{503}{2^{k-3}} - 1, \frac{503}{2^{k-3}}, \frac{503}{2^{k-3}} + 1\]We then have 
 \[503 - 2^{k-3} + 503 > 503 + 2^{k-3}\]
 \[1006 - 2^{k-3} > 503 + 2^{k-3}\]
 \[503 > 2^{k-2}\]
 \[9 > k-2\]
 \[k < 11\]Hence, the first triangle which does not exist in this sequence is $T_{11}$.

 Hence the perimeter is 
 \[\frac{375}{128} + \frac{503}{128} + \frac{631}{128} = \boxed{\textbf{(D)} \frac{1509}{128}}\].

 



\textbf{Solution 2}

Proceeding similarly to the first solution, we have that sides of each triangle are of the form $a, a+1, a+2$ for some number $a$. Also, note that the perimeter of each triangle is half of the previous one. In order for the triangle to not exist, it must not satisfy the triangle inequality, meaning that $a + a + 1 < a+2 \Rightarrow a<1$. Then, the perimeter would be $a + a + 1 + a + 2 = 3a + 3 < 6$. So, to have a proper triangle, we have $\frac{3018}{2^{k}} > 6 \Rightarrow 2^k < 503 \Rightarrow 2^{k} < 512$. The first triangle to not work would have perimeter $\frac{3018}{256} = \frac{1509}{128}$,  thus the answer is   $\boxed{\textbf{(D)} \frac{1509}{128}}$.

 



\textbf{Solution 3}

This problem is a PoP(Power of Point) Bash exercise. What I did to solve this problem was look at recursive perimeters. The first ever perimeter is $P_1 = 2011 + 2012 + 2013$. Next, by PoP, inscribing the circle gives us three new lengths, namely $AD, BE, CF$. Denote $AD$ as $x_1$ and homogeneously the others. Then, $x_1 + x_2 = 2011, x_2 + x_3 = 2012$, and $x_1 + x_3 = 2013$. If we add all these equations up and divide by $2$, we get $x_1 + x_2 + x_3 = \frac{2011 + 2012 + 2013}{2}$. Writing this with $P_1 = 2011 + 2012 + 2013$, we get that our new perimeter, $P_2$, is indeed equal to $\frac{P_1}{2}$. Similarly, by the same concept, we get that $P_3 = \frac{P_1}{4}$ and the pattern keeps going. In general, I found that for each new perimeter $P_n$, $P_n = \frac{P_1}{2^{n-1}}$. Now, substituting in the numerical value of $P_1$, we get that $P_n = \frac{2011 + 2012 + 2013}{2^{n-1}}$. If you keep dividing the numerator and denominator by 2, I got that: $P_n = \frac{1509}{2^{n-3}}$. This representation of a new perimeter in terms of $n$ looks very similar to the option choices, so we're on the right path. Now, all we need to do is find out the last value of $n$ when the sequence stops working. By the Triangular Inequality, we may be able to finish this off. Now, I won't go into depths here, but listing out terms and using the Triangular Inequality gives us:First Sequence:$x_1 < 2012$

 $x_2 < 2013$

 $x_3 < 2011$.The ordered triple for this set of $(x_1, x_2, x_3)$ is $(1005, 1006, 1007)$ if you solve the PoP system of equations we got earlier. If you keep listing out terms, we can come up with the general form, and that is:$x_1 < \frac{503}{2^{n-4}}$

 $x_2 < \frac{503}{2^{n-4}} + 1$

 $x_3 < \frac{503}{2^{n-4}} - 1$.The ordered triple for this set of $(x_1, x_2, x_3)$ is $(\frac{503}{2^{n-4}}, \frac{503}{2^{n-4}} - 1, \frac{503}{2^{n-4}} + 1)$. Now, if we plug in these values of $x_1, x_2$, and $x_3$ into the inequalities, we see that the first two are always satisfied, but the last one is only satisfied when:$2^{n-2} < 503$(You will get this if you simplify the last inequality). The last $n$ for this to be satisfied is when $n = 10$. If we go up to our general representation of $P_n$, we see that $P_n = \frac{1509}{2^{n-3}}$. Plugging in $10$ because this is the last $n$(and also the last triangle), we our final answer of $\frac{1509}{2^{7}} = \boxed{\frac{1509}{128}}$ or $\boxed{D}$.

 



\textbf{Solution 3 (Guess the answer)}

Notice that the perimeter of $T_1$ is $6036$, and the answer choices are $6036(2^{-5}), 6036(2^{-7}), \dots$. So, we assume that the perimeter and each side legnth is roughly cut in half each iteration of $T_n$. We also notice that $2012$ is almost $2^{11}$. The only way that the sequence will end is when the side legnths no longer satisfy the triangle inequality, and since there is a difference of $1$ between the sides of $T_1$, we take a leap of faith and guess $6036(2^{-11})$.~hellohi321

 


\end{enumerate}